% Converted from S4_Template_construction/reviewer_response_package_20261119_claude_v13/RESPONSE_R2_19_EN.md (2026-11-17)
We thank the reviewer and agree. We no longer describe report generation as training-free. The report-realisation stage requires no additional model training or gradient-based fine-tuning, but it depends on template libraries and rules built from training reports. We now describe how these were built and used in the Methods.\par

\textbf{Corpus selection.} For SMU-HMC, 20,400 reports were scored for completeness, and the 400 reports of the test patients were excluded before tokenisation. Of the remaining 20,000 reports, the 7,147 with a completeness score of at least 0.9 were used to build the library. The completeness score is the report-completeness component of ThyClinScore (Methods), computed on the structured version of each report: the weighted fraction of the required information on thyroid measurements, parenchyma, lesions, impression and lymph nodes that is present. The KMVE library was built from the official KMVE training split (1,719 reports); in the revision, the 406 training reports that were near-duplicates of KMVE test reports were excluded, leaving 1,313 reports (our response to Reviewer 1, Comment \hyperref[rev:R1.5]{5}). ZJH-TS reports were generated with the SMU-HMC library.\par

\textbf{Decomposition and de-duplication.} The SMU-HMC reports were divided into five clause categories (thyroid measurement, gland morphology, nodule findings, lymph-node findings and impression). In each clause, case-specific content, such as measurements, locations, echogenicity, vascularity and TI-RADS categories, was replaced by slot placeholders, so that each clause became a candidate template. Identical candidate templates were merged, and the number of clauses mapped to each template was recorded as its frequency. Infrequent templates were then removed: gland and measurement templates outside the 400 and 40 most frequent, which account for 81\% and 99\% of the gland and measurement clauses, respectively; lymph-node templates occurring fewer than 20 times; nodule templates occurring fewer than 3 times; and impression clauses occurring only once. The construction scripts also removed 29 nodule templates whose fixed echogenicity wording conflicted with the echogenicity slot and one measurement template with an incomplete value, and manual review removed 57 gland and measurement templates (see below). In this way, the 1,793 distinct gland templates were reduced to 358, and the 7,478 nodule clauses, which formed 4,656 distinct templates, were reduced to 204 nodule templates. The final SMU-HMC library has 4,774 entries: 400 general sentences (358 gland, 24 measurement and 18 lymph-node sentences), 204 nodule templates, 1,076 impression clauses and 3,094 whole impressions. The KMVE library was rebuilt in the same way during the revision (our response to Reviewer 1, Comment \hyperref[rev:R1.5]{5}): the nodule sentences of these 1,313 training reports were converted into slot templates and merged with their frequencies, giving 82 nodule templates, and the 259 gland sentences were grouped by the gland size, parenchymal echo and parenchymal vascularity they state.\par

\textbf{Manual curation.} Manual review was applied to the 400 most frequent gland templates and the 40 most frequent measurement templates, and removed templates that are not typical of complete routine reports; all removed templates are listed in Appendix A3 (\responsetableref{tab_rep_excluded_templates}). Of the gland templates, 42 were removed: 9 describing postoperative or residual glands and 33 that do not form a complete gland description (9 measurement sentences, 13 describing only glandular vascularity and 11 describing focal lesions). Of the measurement templates, 15 were removed: 7 with a note on an emergency or bedside examination, 4 from postoperative examinations, 3 with missing values and 1 noting that the measurement was unsatisfactory. All other steps were carried out by the construction scripts with the thresholds given above.\par

\textbf{Slots and filling rules.} The SMU-HMC templates contain slots for nodule location, size, composition, echogenicity, margin, shape, aspect ratio, internal echo, echogenic foci and vascularity, and for the TI-RADS category in the impression. A query is built from the structured tool outputs, the best-matching template is retrieved by term matching, and its slots are filled with the corresponding values. Nodules are described by lobe; when both lobes contain nodules, a bilateral template is tried first. The impression is assembled from impression clauses retrieved with the nodule sentences as queries, penalising clauses whose cystic content or calcification disagrees with the nodule sentence, and the TI-RADS category is then filled in. If no clause is retrieved, the most similar whole impression is used; if neither is available, a rule-based impression states the TI-RADS category of each lobe or that no focal lesion was seen. The lymph-node sentence is a fixed sentence reporting enlarged lymph nodes when screening is positive and is otherwise retrieved from the library. When a tool gave no value, the generator used for the reports in the manuscript fell back on defaults: typical nodule features (for example, hypoechoic with a smooth margin), vascularity according to whether a Doppler image was available, and a random value within the typical range for a missing gland dimension; measured gland dimensions were also truncated to a fixed range. Note that our safety rules remove these fallbacks and instead mark the item as not assessed for the physician to complete (our response to Reviewer 1, Comment \hyperref[rev:R1.6]{6}). The KMVE templates contain slots for nodule location, number, echogenicity, margin, shape, echogenic foci and vascularity, which are filled with the same rules; nodule sizes are written as the masked size token of the KMVE reports, and the gland sentence is the most frequent KMVE sentence for the gland size, parenchymal echo and parenchymal vascularity stated by the gland-description tool.\par

\textbf{Centre-specific reporting styles.} Each library reproduces the wording of its source reports, and libraries can be exchanged without changing the image-analysis tools or the evidence schema. SMU-HMC and ZJH-TS reports were generated with the SMU-HMC library, and KMVE reports with the KMVE library, which follows the findings-only format of that dataset. A library for a new centre is built from its reports with the same procedure and added as a separate profile, which is selected when a case is started (Supplementary Fig.~\SuppFigTemplateBank{}). The dependence of the results on the template source is reported in our response to Reviewer 1, Comment \hyperref[rev:R1.5]{5}.\par

\textbf{Clinician assessment.} Two clinicians independently rated, blinded to the system, 270 reports of 90 cases (30 per test set; the ThyroidXAgent report and two baseline reports per case, with the baselines balanced across cases) against the reference report on five dimensions (1--5). Averaged over the 90 cases, the ThyroidXAgent report scored higher than the baseline reports of the same case in factual consistency (2.94 vs 2.46; difference 0.48, 95\% CI 0.24--0.72), readability (4.52 vs 4.08; 0.44, 0.28--0.60) and clinical usability (2.88 vs 2.41; 0.47, 0.22--0.71). By test set, the differences were clear on SMU-HMC (0.68 to 0.79) and ZJH-TS (0.40 to 0.57), mainly because in cases without nodules in the reference report, the baseline reports often described nodules that were not in the reference, whereas the ThyroidXAgent reports rarely did (for factual consistency, +1.43 in the 23 cases without nodules in the reference report, +0.28 in the 33 with one nodule and +0.03 in the 34 with two or more). On KMVE, the differences were small (0.09 to 0.22). We added these results to Results~\SecResReport{}. Agreement between the two clinicians is reported in our response to Reviewer 1, Comment \hyperref[rev:R1.5]{5}; it was low for readability.\par
